\documentclass[final]{beamer}

\usepackage[orientation=portrait,size=a1,scale=1.14]{beamerposter}
\usepackage[utf8]{inputenc}
\usepackage[T1]{fontenc}
\usepackage{lmodern}
\usepackage{amsmath}
\usepackage{booktabs}
\usepackage{fontawesome5}
\usepackage{graphicx}
\usepackage{ragged2e}
\usepackage{xcolor}
\usepackage{tikz}
\usetikzlibrary{arrows.meta,calc,positioning,shapes.geometric}

\graphicspath{{figures/}{../results/cross_run_plots/}}

\definecolor{PosterNavy}{HTML}{233142}
\definecolor{PosterTeal}{HTML}{0E7C7B}
\definecolor{PosterGold}{HTML}{F2A541}
\definecolor{PosterGray}{HTML}{F4F6F8}
\definecolor{PosterInk}{HTML}{172026}

\setbeamercolor{normal text}{fg=PosterInk,bg=white}
\setbeamercolor{headline}{fg=white,bg=PosterNavy}
\setbeamercolor{footline}{fg=white,bg=PosterNavy}
\setbeamercolor{block title}{fg=white,bg=PosterTeal}
\setbeamercolor{block body}{fg=PosterInk,bg=PosterGray}
\setbeamerfont{block title}{size=\large,series=\bfseries}
\setbeamerfont{block body}{size=\normalsize}
\setbeamertemplate{navigation symbols}{}
\setbeamertemplate{caption}[numbered]

\setbeamertemplate{headline}{
  \leavevmode
  \begin{beamercolorbox}[wd=\paperwidth,sep=0.95cm]{headline}
    \begin{columns}[T,totalwidth=\paperwidth]
      \begin{column}{0.76\paperwidth}
        {\Huge\bfseries RAG-NIMBF: A Non-Invasive, Modular Benchmarking Framework for Automated RAG Configuration Analysis\par}
        \vspace{0.45cm}
        {\Large Niclas Hart, Nicolas Weeger \quad Hochschule für Angewandte Wissenschaften Ansbach\par}
      \end{column}
      \begin{column}{0.18\paperwidth}
        \centering
        \vspace{0.2cm}
        \includegraphics[width=0.95\linewidth]{Logo_HS_Ansbach_white.png}
      \end{column}
    \end{columns}
  \end{beamercolorbox}
}

\setbeamertemplate{footline}{
  \leavevmode
  \begin{beamercolorbox}[wd=\paperwidth,sep=0.42cm]{footline}
    \large n.hart18293@hs-ansbach.de \hfill ARC, 03.07.2026
  \end{beamercolorbox}
}

\newcommand{\metric}[2]{\textbf{\textcolor{PosterTeal}{#1}} #2}
\newcommand{\todo}[1]{\textcolor{PosterGold}{\textbf{TODO:} #1}}
\newcommand{\blockheading}[2]{\faIcon{#1}\hspace{0.45em}#2}

\title{RAG-NIMBF: A Non-Invasive, Modular Benchmarking Framework for Automated RAG Configuration Analysis}
\author{Niclas Hart}
\institute{Hochschule fuer Angewandte Wissenschaften Ansbach}
\date{July 2026}

\begin{document}
\begin{frame}[t]

\begin{columns}[T,totalwidth=\textwidth]

\begin{column}{0.31\textwidth}

\begin{block}{\blockheading{question-circle}{Problem}}
\justifying
Retrieval-augmented generation (RAG) systems expose a large configuration
space. Chunking, overlap, embedding model, vector database, retrieval strategy,
reranking, prompt template, generator model, and evaluator choice can all move
quality and runtime in different directions.

\vspace{0.4cm}
\textbf{Research goal:} make RAG configuration choices measurable,
repeatable, and comparable without rewriting the target pipeline for each
experiment.
\end{block}

\begin{block}{\blockheading{puzzle-piece}{Framework Contribution}}
\begin{itemize}\setlength\itemsep{0.05cm}
  \item Environment-driven benchmark grids for local and OpenAI-compatible
        model providers.
  \item Dataset adapters normalize QA corpora into one benchmark record shape.
  \item Reusable retrieval core supports Chroma and LanceDB vector stores,
        similarity search, MMR, optional HyDE, and cross-encoder reranking.
  \item Generation layer captures answer text, time to first token, throughput,
        and optional GPU statistics.
  \item Reporting exports JSON, CSV, Markdown, static plots, MLflow runs, and
        optional tracing artifacts.
\end{itemize}
\end{block}

\begin{block}{\blockheading{project-diagram}{Pipeline}}
\centering
\vspace{0.35cm}
\begin{tikzpicture}[
  node distance=0.95cm,
  stage/.style={
    draw=PosterTeal,
    fill=white,
    rounded corners=6pt,
    line width=1.55pt,
    text width=0.88\linewidth,
    align=center,
    minimum height=1.30cm,
    inner sep=0.18cm,
    font=\normalsize\bfseries
  },
  side/.style={
    draw=PosterGold,
    fill=PosterGold!12,
    rounded corners=6pt,
    line width=1.35pt,
    text width=0.40\linewidth,
    align=center,
    minimum height=1.35cm,
    inner sep=0.28cm,
    font=\normalsize\bfseries
  },
  flow/.style={-Latex, line width=1.75pt, draw=PosterNavy, shorten >=4pt, shorten <=4pt}
]
\node[side] (config) {Configuration grid\\models, chunks, prompts};
\node[side, right=1.0cm of config] (dataset) {Dataset adapter\\QA record shape};
\coordinate (inputsouth) at ($(config.south)!0.5!(dataset.south)$);
\node[stage, anchor=north] (chunk) at ($(inputsouth)+(0,-1.35cm)$) {Chunking + vector-store cache\\reusable index for each corpus/configuration};
\node[stage, below=0.60cm of chunk] (retrieve) {Retrieval, optional reranking, prompt assembly\\similarity/MMR contexts become model input};
\node[stage, below=0.60cm of retrieve] (generate) {Generation\\answer text, TTFT, throughput, GPU stats};
\node[stage, below=0.60cm of generate] (evaluate) {RAGAS + custom IR/NLG metrics\\faithfulness, similarity, retrieval quality, runtime};
\node[stage, below=0.60cm of evaluate] (report) {Artifacts\\JSON, CSV, Markdown, plots, MLflow, traces};

\draw[flow] (config) -- (chunk);
\draw[flow] (dataset) -- (chunk);
\draw[flow] (chunk) -- (retrieve);
\draw[flow] (retrieve) -- (generate);
\draw[flow] (generate) -- (evaluate);
\draw[flow] (evaluate) -- (report);
\end{tikzpicture}
\end{block}

\begin{block}{\blockheading{sliders-h}{Experiment Space}}
\begin{tabular}{ll}
\toprule
Dimension & Current values in project \\
\midrule
Chunking & recursive, semantic, text, token, markdown \\
Chunk size & 500, 1000, 1500 and variants \\
Overlap & 100, 200, 350 and variants \\
Retrieval & similarity, MMR \\
Reranking & none, MiniLM cross-encoder \\
Embedding & \texttt{nomic-embed-text:latest} snapshot \\
Datasets & SQuAD / WikiQA-style QA adapters \\
Models & Qwen3, Qwen3.5, Qwen3.6, GPT-OSS, GPT-4o-mini, DeepSeek \\
\bottomrule
\end{tabular}
\end{block}

\begin{block}{\blockheading{redo}{Reproducibility Controls}}
\begin{itemize}\setlength\itemsep{0.04cm}
  \item \metric{Config grid:}{environment-defined combinations generate every run.}
  \item \metric{Run folders:}{each benchmark writes to \texttt{results/runN/}.}
  \item \metric{Vector cache:}{corpus fingerprints avoid accidental re-embedding.}
  \item \metric{Stage timings:}{indexing, retrieval, generation, and evaluation are timed.}
  \item \metric{Artifact trail:}{JSON, CSV, plots, MLflow, and traces preserve evidence.}
\end{itemize}
\end{block}

\begin{block}{\blockheading{lightbulb}{Observed Answers}}
\begin{itemize}\setlength\itemsep{0.04cm}
  \item \metric{Prompting:}{across \(N \geq 100\) rows, detailed prompts raise mean faithfulness (0.908, \(n=125\)) over concise (0.838, \(n=132\)), while BERT-F1 stays nearly unchanged (0.873 vs. 0.871).}
  \item \metric{Retrieval:}{similarity is stronger for faithfulness (0.886, \(n=146\)) vs. MMR (0.854, \(n=111\)), while MMR improves nDCG@5 (0.779 vs. 0.686) and recall@5 (0.826 vs. 0.559).}
  \item \metric{Chunking:}{recursive and semantic chunking are close on mean faithfulness (0.875, \(n=158\) vs. 0.867, \(n=99\)), while semantic improves mean nDCG@5 (0.747 vs. 0.713).}
  \item \metric{Winner:}{no single configuration dominates: faithfulness, BERT-F1, nDCG@5, recall@5, METEOR, and runtime select different runs.}
\end{itemize}
\end{block}


\end{column}

\begin{column}{0.36\textwidth}

\begin{block}{\blockheading{search}{Evaluation Lenses}}
\begin{itemize}
  \item \metric{LLM-judged quality:}{faithfulness, context recall,
        semantic similarity, and related RAGAS fields where available.}
  \item \metric{Retrieval quality:}{hit@k, nDCG@k, recall@k, context relevance,
        with heuristic or gold-document modes.}
  \item \metric{Answer similarity:}{ROUGE-L, BLEU, METEOR, and optional
        BERTScore precision, recall, F1.}
  \item \metric{Performance:}{total runtime, stage timings, time to first
        token, tokens per second, and GPU usage when available.}
\end{itemize}
\end{block}


\begin{block}{\blockheading{comment-dots}{Prompt Templates}}
\setlength{\tabcolsep}{0.35cm}
\renewcommand{\arraystretch}{1.08}
\begin{tabular}{@{}p{0.48\linewidth}p{0.48\linewidth}@{}}
\textbf{\textcolor{PosterTeal}{Concise Answer Prompt}} &
\textbf{\textcolor{PosterTeal}{Detailed Answer Prompt}} \\
\midrule
\raggedright
\scriptsize\ttfamily
Answer the question using ONLY the provided context.\\
RULES:\\
- Give a short, direct, factual answer -- one sentence maximum.\\
- Do NOT explain your reasoning, add context, or use markdown formatting.\\
- Do NOT say 'Based on the context' or similar preamble.\\
- If the context does not contain the answer, say 'I cannot answer from the provided context.'\\
Examples:\\
Q: What circuit court is Maryland? \(\rightarrow\) The Circuit Courts of Maryland are the state trial courts of general jurisdiction.\\
Q: How many members are in the House of Representatives? \(\rightarrow\) The total number of voting representatives is fixed at 435.\\
Q: Who wrote the music for Star Wars? \(\rightarrow\) John Williams.
&
\raggedright
\scriptsize\ttfamily
Answer the question based only on the provided context. Provide a clear, complete answer in full sentences. Include relevant calculations or reasoning steps when applicable. Do not start your answer with phrases like 'Based on the provided context' or 'Based on the given context'. Answer directly.
\end{tabular}

\end{block}

\begin{block}{\blockheading{chart-bar}{Current Results Snapshot}}
\justifying
The current result corpus contains 274 JSON configuration rows across 86 benchmark reports, with 257 rows at \(N \geq 100\).
The manual evaluation matrix tracks configuration IDs up to 222 and currently marks 212 full \(N=100\) rows as tested.
The table below highlights the strongest observed configurations for three evaluation lenses, demonstrating that different configurations can excel in different metrics,
underscoring the importance of multi-metric evaluation in RAG benchmarking.
\vspace{0.25cm}
{\scriptsize
\setlength{\tabcolsep}{0.12cm}
\renewcommand{\arraystretch}{1.02}
\noindent\resizebox{0.98\linewidth}{!}{%
\begin{tabular}{p{0.42\linewidth}cccc}
\toprule
Observed strongest configuration & Faith. & BERT-F1 & nDCG@5 & N \\
\midrule
Qwen3.5-397B, detailed, recursive 500/100, similarity &
0.959 & 0.835 & 0.702 & 300 \\
Qwen3-32B-AWQ, concise, semantic, MMR &
0.871 & 0.969 & 1.000 & 100 \\
deepseek-v4-pro, concise, recursive 1000/100, similarity &
0.845 & 0.919 & 0.675 & 100 \\
\bottomrule
\end{tabular}%
}
}

\vspace{0.25cm}
\textbf{Interpretation:} the best run depends on the metric. Strong
faithfulness, strong retrieval ranking, and strong lexical/semantic answer
similarity do not always select the same configuration.
\end{block}

\begin{block}{\blockheading{sort-amount-up}{Quality vs. Speed}}
\centering
\includegraphics[width=0.92\linewidth]{paper_quality_vs_speed.png}

\vspace{0.1cm}
\justifying
{\footnotesize Faithfulness vs. throughput; dashed line = Pareto frontier.
No single config dominates — frontier trades quality for speed across
LLMs.}
\end{block}

\begin{block}{\blockheading{robot}{Prompt and Model Effects}}
\centering
\includegraphics[width=0.88\linewidth]{paper_llm_metric_heatmap.png}
\end{block}

\end{column}

\begin{column}{0.31\textwidth}

\begin{block}{\blockheading{balance-scale}{Retrieval Trade-Offs}}
\centering
\includegraphics[width=0.88\linewidth]{paper_grouped_ragas_context_recall.png}

\vspace{0.2cm}
\justifying
MMR can improve diversity-sensitive retrieval metrics in some runs, while
plain similarity can preserve higher faithfulness in others. Reranking helps in
selected settings, but adds latency and should be judged jointly with answer
quality.
\end{block}

\begin{block}{\blockheading{tachometer-alt}{LLM Runtime Baseline}}
\centering
{\footnotesize
\setlength{\tabcolsep}{0.10cm}
\renewcommand{\arraystretch}{1.10}
\begin{tabular*}{\linewidth}{@{\extracolsep{\fill}}p{0.60\linewidth}rr@{}}
\toprule
Model / backend & Latency & Gen tok/s \\
\midrule
Qwen3.5-397B / Cloud & 7.82s & 130.9 \\
Qwen3.5-Think / Cloud & 7.84s & 130.6 \\
Qwen3.6-27B / RTX5090 vLLM & 8.67s & 118.1 \\
GPT-4o-mini / Cloud & 10.12s & 101.1 \\
Qwen3-32B / RTX5090 vLLM & 17.36s & 59.0 \\
gpt-oss-20B / Spark Ollama & 18.13s & 56.5 \\
\bottomrule
\end{tabular*}
}

\vspace{0.22cm}
\justifying
{\footnotesize Independent LLM performance baseline from
\texttt{LLM\_Performance\_Tests-main}: sequential single-call runs with
\(N=1\), max output 1024 tokens. Throughput is generated output tokens per
second. TTFT is omitted from the table because two traces are non-streamed;
output-token count is a fixed benchmark cap rather than a differentiating
metric.}
\end{block}

\begin{block}{\blockheading{exclamation-triangle}{Limitations}}
\normalsize
\begin{itemize}\setlength\itemsep{0.05cm}
  \item Current results are local, hardware-dependent, and partly duplicated by
        reruns.
  \item SQuAD/WikiQA-style corpora may not transfer to domain RAG workloads.
  \item LLM evaluator reliability, stochastic generation, and provider changes
        remain validity threats.
  \item Manual matrix entries must be reconciled with generated result files
        before publication.
\end{itemize}
\end{block}

\begin{block}{\blockheading{clipboard-check}{Takeaways for Practitioners}}
\begin{enumerate}
  \item Benchmark complete RAG configurations, not isolated components.
  \item Report multiple metrics because each metric rewards different behavior.
  \item Store machine-readable artifacts for every run.
  \item Treat prompt template and retrieval strategy as first-class tuning
        variables.
  \item Keep limitations visible when using LLM-judged metrics.
\end{enumerate}
\end{block}

\begin{block}{\blockheading{road}{Future Work}}
\begin{itemize}\setlength\itemsep{0.05cm}
  \item \metric{Datasets:}{extend beyond SQuAD/WikiQA to domain corpora
        (legal, medical, enterprise KB).}
  \item \metric{Scale:}{raise \(N\) per config and add significance tests
        so single-run noise is visible.}
  \item \metric{Retrieval ablation:}{full HyDE and cross-encoder sweep,
        plus hybrid lexical+dense retrieval.}
  \item \metric{Backends:} compare Chroma vs. LanceDB on index time,
        memory, and query latency at fixed corpus size.
\end{itemize}
\end{block}

\begin{block}{\blockheading{chart-line}{Statistical Robustness}}
\centering
\includegraphics[width=0.66\linewidth]{ci_faithfulness_by_prompt.png}

\vspace{0.06cm}
{\footnotesize Bootstrap 95\% CI, paired Wilcoxon \(p<0.001\).}
\end{block}

\begin{block}{\blockheading{book}{References}}
\begin{columns}[T,totalwidth=\linewidth]
  \begin{column}{0.68\linewidth}
    \scriptsize
    \begin{enumerate}\setlength\itemsep{0.0cm}
      \item Lewis et al. \emph{RAG for Knowledge-Intensive NLP Tasks.} NeurIPS 2020.
      \item Es et al. \emph{RAGAS: Automated Evaluation of RAG.} arXiv:2309.15217, 2023.
      \item Rajpurkar et al. \emph{SQuAD: 100{,}000+ Questions for Machine Comprehension.} EMNLP 2016.
      \item Lin. \emph{ROUGE: A Package for Automatic Evaluation of Summaries.} 2004.
      \item Zhang et al. \emph{BERTScore: Evaluating Text Generation with BERT.} ICLR 2020.
    \end{enumerate}
  \end{column}
  \begin{column}{0.24\linewidth}
    \centering
    \href{https://github.com/niclashart/Benchmarking-Framework}{\includegraphics[width=0.95\linewidth]{qr-code.png}}\\[-0.05cm]
    {\scriptsize scan for code}
  \end{column}
\end{columns}
\end{block}


\end{column}

\end{columns}

\end{frame}
\end{document}
